\subsection{Поворот и проецирование}

Перед отображением все вершины последовательно поворачиваются вокруг осей X, Y и Z. Центр вращения совпадает с центром фигуры. Углы задаются ползунками интерфейса или изменяются перетаскиванием мыши. Поворот сохраняет форму, длины рёбер и расстояния между вершинами.

Для вектора-столбца координат порядок преобразований записывается как
\begin{equation}
  p'=R_z(\gamma)R_y(\beta)R_x(\alpha)p,
  \label{eq:rotation}
\end{equation}
где \(\alpha,\beta,\gamma\) — углы поворота вокруг осей X, Y и Z.
Матрица справа применяется первой. Например, поворот вокруг X имеет вид
\begin{equation}
  R_x(\alpha)=
  \begin{pmatrix}
    1 & 0 & 0 \\
    0 & \cos\alpha & -\sin\alpha \\
    0 & \sin\alpha & \cos\alpha
  \end{pmatrix}.
  \label{eq:rotation-x}
\end{equation}
В коде те же операции выполняются непосредственно над компонентами вектора,
без хранения отдельных матриц.

Камера расположена на положительной полуоси Z и направлена к центру фигуры. Плоскость проекции проходит через начало координат. Положение камеры фиксировано, а пользователь изменяет ориентацию объекта. Нормали граней поворачиваются вместе с вершинами, чтобы проверка видимости и освещение соответствовали текущему положению модели.

Ортографическая проекция переносит горизонтальную и вертикальную составляющие положения вершины на плоскость экрана без уменьшения удалённых частей фигуры. Размер изображения не зависит от глубины точки. Такой режим удобен для изучения формы и сравнения пропорций.

Для точки \((x,y,z)\) после поворота ортографическая проекция задаётся выражениями
\begin{equation}
  x_{\mathrm{pr}}=x, \qquad y_{\mathrm{pr}}=y.
  \label{eq:orthographic}
\end{equation}

При перспективной проекции учитывается расстояние от вершины до камеры. Ближние части объекта выглядят крупнее, а дальние — меньше. Это создаёт ощущение глубины. В программе проекция вычисляется непосредственно для каждой вершины на центральном процессоре; отдельные матрицы проекции в шейдер не передаются.

При расположении камеры в точке \((0,0,d)\) и плоскости проекции \(z=0\)
используются формулы
\begin{equation}
  x_{\mathrm{pr}}=\frac{dx}{d-z}, \qquad
  y_{\mathrm{pr}}=\frac{dy}{d-z}, \qquad d=6.
  \label{eq:perspective}
\end{equation}
Все вершины при допустимых поворотах находятся перед камерой, поэтому
\(d-z>0\). В экранной системе результат умножается на общий масштаб,
смещается к центру области рисования, а вертикальная координата меняет знак.
